%=====================================================================
% MONTAGEM PREDIAL - SEGURANCA, PLANEJAMENTO E FERRAMENTAS
%=====================================================================

\section{Planejamento seguro, ferramentas e materiais}

\begin{frame}{Montagem começa antes de abrir a primeira caixa}
\centering
\begin{tikzpicture}[
  b/.style={draw=corPrim,fill=corDef!7,rounded corners,minimum width=2.15cm,minimum height=0.78cm,align=center,font=\scriptsize},
  a/.style={-{Stealth},thick,corPrim}]
\node[b](proj)at(-5.0,0){Projeto\\liberado};
\node[b](risco)at(-2.5,0){Análise de risco\\e procedimento};
\node[b](mat)at(0,0){Materiais e\\instrumentos};
\node[b](exec)at(2.5,0){Execução\\controlada};
\node[b](ensaio)at(5.0,0){Inspeção, ensaio\\e registro};
\draw[a](proj)--(risco);\draw[a](risco)--(mat);\draw[a](mat)--(exec);\draw[a](exec)--(ensaio);
\end{tikzpicture}
\vspace{5mm}
\begin{caixaAt}
Este material é didático. Intervenções em instalações elétricas exigem trabalhadores qualificados, capacitados e autorizados conforme as responsabilidades e condições aplicáveis.
\end{caixaAt}
\end{frame}

\begin{frame}{Responsabilidades precisam estar definidas}
\small
\begin{tabularx}{\textwidth}{>{\bfseries}p{2.9cm}X X}
\toprule
Papel & Antes do serviço & Durante e depois \\
\midrule
Responsável técnico & projeto, especificações e critérios de aceitação & decisões técnicas, alterações e documentação \\
Supervisor & equipe, recursos, análise de risco e sequência & controle da execução e liberação das etapas \\
Executor autorizado & conferir tarefa, ferramentas e condição segura & executar, comunicar desvios e registrar resultados \\
Responsável por ensaios & método, instrumento e limites & medir, interpretar e emitir registro \\
Cliente/operação & acesso, cargas e restrições de uso & aceite, operação e manutenção futura \\
\bottomrule
\end{tabularx}
\end{frame}

\begin{frame}{Ordem de serviço: transformar projeto em tarefa controlada}
\small
\begin{enumerate}
\item identificar local, circuito, origem e limites físicos da atividade;
\item anexar desenho, revisão e lista de materiais aplicáveis;
\item registrar perigos elétricos e não elétricos;
\item definir desenergização, impedimento de reenergização e liberação;
\item atribuir responsáveis, equipe e comunicação;
\item listar instrumentos, EPIs, EPCs e calibração necessária;
\item estabelecer inspeções intermediárias e critérios de aceite;
\item fechar alterações, ensaios, fotos e as built.
\end{enumerate}
\begin{caixaObs}
Uma ordem de serviço genérica não controla risco nem qualidade. A tarefa precisa ter origem, destino, revisão e critério de término.
\end{caixaObs}
\end{frame}

\begin{frame}{Análise de risco: eletricidade não é o único perigo}
\small
\begin{tabularx}{\textwidth}{>{\bfseries}p{2.9cm}X X}
\toprule
Perigo & Exposição típica & Controle planejado \\
\midrule
Choque/arco & quadro, condutor existente, retorno e geração & desenergização, teste, bloqueio, barreiras \\
Altura & luminárias, fachadas e entrada aérea & acesso adequado, NR-35 quando aplicável \\
Perfuração & parede com tubulação ou circuito oculto & compatibilização, detecção e delimitação \\
Poeira/projeção & cortes, rasgos e furações & proteção coletiva, ocular e respiratória \\
Esforço/queda & bobinas, escadas, ferramentas e entulho & organização, movimentação e área isolada \\
Ambiente & umidade, calor, chuva e espaço reduzido & suspender, adaptar método ou controlar condição \\
\bottomrule
\end{tabularx}
\end{frame}

\begin{frame}{Ferramentas: escolher pelo processo, não pelo hábito}
\small
\begin{tabularx}{\textwidth}{>{\bfseries}p{2.6cm}X X}
\toprule
Etapa & Ferramentas principais & Controle \\
\midrule
Marcação & trena, nível, esquadro, linha e detector & cotas e interferências \\
Infraestrutura & furadeira, serra-copo, curvador, guia e escareador & diâmetro, raio e acabamento \\
Condutores & guia, lubrificante compatível, desenrolador e puxador & tração e torção \\
Preparação & decapador, alicate de corte e desencapador de cabo & não ferir a alma \\
Terminação & prensa-terminal e torquímetro/chave de torque & matriz e torque corretos \\
Identificação & impressora/etiqueta e marcador permanente & padrão único e legível \\
\bottomrule
\end{tabularx}
\end{frame}

\begin{frame}{Instrumentos: cada um responde a uma pergunta}
\small
\begin{tabularx}{\textwidth}{>{\bfseries}p{2.9cm}X X}
\toprule
Instrumento & Pergunta & Limite de uso \\
\midrule
Detector bipolar & há tensão e qual a faixa? & verificar antes/depois em fonte conhecida \\
Multímetro TRMS & tensão, resistência, frequência? & categoria, ponta e fusível adequados \\
Alicate amperímetro & corrente sem abrir o circuito? & abraçar o condutor correto, não o cabo completo \\
Megômetro & isolação suporta tensão de ensaio? & desconectar eletrônica sensível \\
Terrômetro & resistência do sistema/eletrodo? & método e geometria compatíveis \\
Testador de DR & corrente/tempo de atuação? & interpretar conforme dispositivo e circuito \\
Analisador de rede & harmônicas, potência e eventos? & configuração, janela e referência de tensão \\
\bottomrule
\end{tabularx}
\end{frame}

\begin{frame}{Ferramenta isolada não autoriza trabalho energizado}
\small
\begin{columns}[T,onlytextwidth]
\column{0.49\textwidth}
\begin{caixaAt}[title=Erro perigoso]
Confundir cabo isolado, cabo de ferramenta revestido ou luva comum com proteção certificada para o risco elétrico.
\end{caixaAt}
\column{0.49\textwidth}
\begin{caixaNorma}[title=Uso correto]
\begin{itemize}
\item ferramenta compatível com a tarefa e tensão;
\item inspeção visual antes do uso;
\item armazenamento que preserve a isolação;
\item retirada imediata quando houver dano;
\item procedimento de trabalho prevalece sobre o acessório.
\end{itemize}
\end{caixaNorma}
\end{columns}
\begin{caixaObs}
Trabalho energizado é exceção tecnicamente justificada, com análise e controles específicos. A regra de montagem predial é preparar e executar desenergizado.
\end{caixaObs}
\end{frame}

\begin{frame}{EPI e EPC: barreiras complementares}
\small
\begin{tabularx}{\textwidth}{>{\bfseries}p{2.7cm}X X}
\toprule
Grupo & Exemplos & Verificar \\
\midrule
Proteção coletiva & bloqueio, barreira, sinalização, cobertura e aterramento temporário quando aplicável & instalação, estabilidade e responsabilidade \\
Cabeça/face & capacete e proteção facial compatível & classe, estado e ajuste \\
Mãos & luvas adequadas ao risco e cobertura mecânica quando necessária & ensaio, validade e perfuração \\
Corpo & vestimenta compatível com o risco térmico e ambiente & fechamento, conservação e contaminação \\
Olhos/respiração & óculos, viseira e proteção respiratória & poeira, projeção e ventilação \\
Altura & sistema de acesso e proteção contra queda & inspeção e ancoragem \\
\bottomrule
\end{tabularx}
\begin{caixaAt}
EPI não corrige circuito energizado sem necessidade, bloqueio ausente ou procedimento incompleto.
\end{caixaAt}
\end{frame}

\begin{frame}{Recebimento: rejeitar material antes de instalar}
\small
\begin{columns}[T,onlytextwidth]
\column{0.49\textwidth}
\begin{itemize}
\item conferir fabricante, modelo e certificação aplicável;
\item comparar seção, tensão, temperatura e classe do cabo;
\item verificar curva, polos, corrente e capacidade do disjuntor;
\item confirmar tipo e sensibilidade do DR;
\item conferir $U_c$, $U_p$, $I_n/I_{max}$ do DPS;
\item inspecionar trincas, oxidação, umidade e lacres.
\end{itemize}
\column{0.49\textwidth}
\begin{caixaProjeto}[title=Rastreabilidade]
Registrar:
\begin{itemize}
\item lote e nota de recebimento;
\item ficha técnica usada;
\item quantidade aceita/rejeitada;
\item local de armazenamento;
\item substituições aprovadas.
\end{itemize}
\end{caixaProjeto}
\end{columns}
\end{frame}

\begin{frame}{Organização do posto reduz falha de montagem}
\small
\begin{columns}[T,onlytextwidth]
\column{0.52\textwidth}
\begin{tikzpicture}[scale=0.82]
\draw[fill=black!3] (0,0) rectangle (6.2,3.8);
\draw[fill=corDef!8,draw=corPrim] (0.35,2.45) rectangle (2.25,3.35);
\node[font=\scriptsize,align=center] at (1.3,2.9){componentes\\conferidos};
\draw[fill=corNorma!8,draw=corNorma] (2.55,2.45) rectangle (4.45,3.35);
\node[font=\scriptsize,align=center] at (3.5,2.9){ferramentas\\inspecionadas};
\draw[fill=corAccent!10,draw=corAccent] (4.75,2.45) rectangle (5.85,3.35);
\node[font=\scriptsize,align=center] at (5.3,2.9){resíduos};
\draw[very thick] (0.45,0.75)--(5.75,0.75);
\node[font=\scriptsize] at (3.1,1.15){bancada limpa e iluminada};
\node[font=\tiny] at (3.1,0.35){documento e circuito identificados};
\end{tikzpicture}
\column{0.46\textwidth}
\begin{itemize}
\item separar materiais por circuito;
\item manter pontas e ferramentas protegidas;
\item evitar sobras dentro de caixas e quadros;
\item controlar parafusos, terminais e etiquetas;
\item proteger equipamentos contra poeira e impacto;
\item registrar etapa concluída antes de fechar.
\end{itemize}
\end{columns}
\end{frame}

\begin{frame}{Liberação da etapa de montagem}
\small
\begin{enumerate}
\item projeto e revisão corretos estão disponíveis no local;
\item circuito e fontes possíveis foram identificados;
\item desenergização, bloqueio e teste foram concluídos;
\item equipe, responsabilidades e comunicação estão definidos;
\item ferramentas, instrumentos e EPIs/EPCs foram inspecionados;
\item material confere com a especificação;
\item ambiente permite execução segura e limpa;
\item critérios de inspeção e registro estão definidos.
\end{enumerate}
\begin{caixaAt}
Qualquer resposta negativa interrompe a liberação. Corrigir a condição antes de iniciar a execução.
\end{caixaAt}
\end{frame}
